%% ch06-redis-schema.tex — Redis Schema
%% JA4proxy Technical Reference Manual

\chapter{Redis Schema}
\index{Redis!schema}
\label{ch:redis}

\begin{chapterquote}
Redis is the shared state layer. Every \ja\ instance reads and writes the same
Redis instance. All bans, rate windows, beaconing timestamps, JA4 lists, and
enrichment caches live here. The local in-process cache is a performance overlay —
Redis is the source of truth.
\end{chapterquote}

\newthought{Redis data structure selection} is not arbitrary. Each use case has a
specific rationale — wrong data structures produce incorrect sliding windows,
lost updates, or excessive memory usage. This chapter documents every key pattern,
its data structure, and the rationale for that choice.

\section{Data Structure Rationale}
\index{Redis!data structures}

\begin{fullwidth}
\begin{table}[h]
\caption{Redis data structure selection rationale}
\label{tab:redis-rationale}
\begin{tabular}{@{}llp{4cm}p{4.5cm}@{}}
\toprule
\textbf{Use case} & \textbf{Structure} & \textbf{Why this structure} & \textbf{Alternative considered} \\
\midrule
JA4 blacklist/whitelist & SET & O(1) SISMEMBER; small and static & Hash — unnecessary overhead \\
IP bans & String + TTL & Per-key TTL required; atomic SETEX & SET — no per-key TTL \\
Sliding window rate limiting & Sorted Set + Lua (EVALSHA) & True sliding window; no boundary errors & INCR + EXPIRE — fixed window only \\
Beaconing timestamps & Sorted Set (score=timestamp) & ZRANGEBYSCORE for time windows; ordered & List — no time range queries \\
Session resumption ratio & Hash (total, resumed) & Atomic HINCRBY; two fields one key & Two strings — two round-trips \\
Concurrent connection count & String (INCR/DECR) & Simple atomic counter & Hash — unnecessary \\
Return visitor tracking & Hash (first\_seen, total, allowed, blocked) & Multi-field HINCRBY; atomic & Multiple keys — no atomicity \\
Unique IP per CIDR & HyperLogLog (PFADD/PFCOUNT) & O(1), 12KB/key, 0.81\% error & Hash of IPs — unbounded memory \\
Enrichment dedup & Bloom filter (BF.ADD/BF.EXISTS) & False positives acceptable; O(1) & SET — exact but uses more memory \\
CIDR matching & In-process trie (NOT Redis) & No CIDR primitive in Redis; O(log n) trie & Redis SCAN — too slow \\
GeoIP/ASN lookup & In-process mmap (NOT Redis) & MaxMind designed for mmap & Redis hash — defeats mmap optimisation \\
Cross-instance events & Stream (XADD/XREADGROUP) & Persistent, replayable; consumer groups & Pub/Sub — ephemeral, no replay \\
Analytics findings & String + TTL & Read by proxy; simple JSON blob & Stream — adds complexity \\
\bottomrule
\end{tabular}
\end{table}
\end{fullwidth}

\begin{dangerbox}[Never Use Redis for CIDR Matching]
CIDR matching in Redis requires iterating all keys matching a pattern or maintaining
a sorted set of prefixes. Both approaches are O(n) for large blocklists. The in-process
trie (pytricia for Python, prefix trie for Go) provides O(log n) lookup and is the
only correct approach. All CIDR matching is done in-process; Redis stores the list
for sharing across instances and reloading after restart.
\end{dangerbox}

\section{JA4 Fingerprint Lists}
\index{Redis!ja4:blacklist}
\index{Redis!ja4:whitelist}

\subsection{ja4:blacklist}

\begin{fullwidth}
\begin{table}[h]
\caption{\code{ja4:blacklist} key reference}
\label{tab:redis-ja4-blacklist}
\begin{tabular}{@{}ll@{}}
\toprule
\textbf{Key} & \code{ja4:blacklist} \\
\textbf{Type} & SET \\
\textbf{TTL} & None (persistent) \\
\textbf{Writer} & Config loader at startup; Management UI \\
\textbf{Reader} & Pipeline layer 2 (JA4 blacklist bypass) \\
\textbf{Example members} & \code{t13d190900\_9dc949149365\_97f8aa674fd9} \\
\textbf{Operations} & \code{SADD}, \code{SREM}, \code{SISMEMBER}, \code{SMEMBERS} \\
\bottomrule
\end{tabular}
\end{table}
\end{fullwidth}

At startup, the config loader reads \configkey{security.blacklist} from
\code{proxy.yml} and executes a Redis pipeline to \code{SADD} all entries to
\code{ja4:blacklist}. On hot-reload, entries removed from config are \code{SREM}'d.
An in-process set mirrors the Redis SET for zero-latency lookup on the hot path.

\subsection{ja4:whitelist}

\begin{fullwidth}
\begin{table}[h]
\caption{\code{ja4:whitelist} key reference}
\label{tab:redis-ja4-whitelist}
\begin{tabular}{@{}ll@{}}
\toprule
\textbf{Key} & \code{ja4:whitelist} \\
\textbf{Type} & SET \\
\textbf{TTL} & None (persistent) \\
\textbf{Writer} & Config loader at startup; Management UI \\
\textbf{Reader} & Pipeline layer 1b (JA4 whitelist bypass) \\
\textbf{Example members} & \code{t13d1516h2\_8daaf6152771\_02713d6af862} \\
\bottomrule
\end{tabular}
\end{table}
\end{fullwidth}

\section{IP Bans}
\index{Redis!ban}

\subsection{ban:\{ip\}}

\begin{fullwidth}
\begin{table}[h]
\caption{\code{ban:\{ip\}} key reference}
\label{tab:redis-ban}
\begin{tabular}{@{}ll@{}}
\toprule
\textbf{Key pattern} & \code{ban:\{ip\}} where \code{\{ip\}} is the canonical IP string \\
\textbf{Type} & String \\
\textbf{Value} & JSON: \code{\{"banned\_at": 1712232896, "reason": "rate\_limit\_by\_ip", "ttl": 300\}} \\
\textbf{TTL} & \code{ban\_duration} seconds (escalates on re-ban) \\
\textbf{Writer} & Action decider (ban action) \\
\textbf{Reader} & Pipeline layer 0 (early ban check); local LRU cache \\
\textbf{Examples} & \code{ban:185.220.101.1}, \code{ban:2a06:98c1:3120::3} \\
\textbf{Operations} & \code{SET key value EX ttl}, \code{GET}, \code{DEL} \\
\bottomrule
\end{tabular}
\end{table}
\end{fullwidth}

Ban TTL follows an escalation schedule: initial TTL is \configkey{ban\_duration}
(default 300s); subsequent bans double the TTL up to a maximum of 604800s (7 days).
The current escalation level is stored in the JSON value's \code{escalation} field.

\section{Rate Limiting Windows}
\index{Redis!rate limiting}
\index{sliding window}

\subsection{rate:\{strategy\}:\{key\}}

\begin{fullwidth}
\begin{table}[h]
\caption{\code{rate:\{strategy\}:\{key\}} key reference}
\label{tab:redis-rate}
\begin{tabular}{@{}ll@{}}
\toprule
\textbf{Key pattern} & \code{rate:\{strategy\}:\{key\}} \\
\textbf{Type} & Sorted Set \\
\textbf{Score} & Unix timestamp in milliseconds \\
\textbf{Member} & Connection UUID (ensures uniqueness) \\
\textbf{TTL} & Managed by ZREMRANGEBYSCORE; no explicit key TTL \\
\textbf{Writer} & Pipeline rate-limiting layer \\
\textbf{Reader} & Pipeline rate-limiting layer (ZCOUNT) \\
\textbf{Example keys} & \code{rate:by\_ip:185.220.101.1} \\
 & \code{rate:by\_ja4:t13d190900\_...} \\
 & \code{rate:by\_ip\_ja4\_pair:185.220.101.1:t13d190900\_...} \\
\textbf{Operations} & Lua EVALSHA for atomic ZADD + ZREMRANGEBYSCORE + ZCOUNT \\
\bottomrule
\end{tabular}
\end{table}
\end{fullwidth}

The sliding window is implemented as a Lua script that atomically:
\begin{enumerate}
  \item Removes entries older than the window size (ZREMRANGEBYSCORE)
  \item Adds the current timestamp (ZADD)
  \item Returns the current window count (ZCOUNT)
\end{enumerate}

Using Lua (EVALSHA) ensures the three operations are atomic from Redis's perspective —
no race condition is possible between the remove and count steps.

\section{Beaconing Data}
\index{Redis!beaconing}

\subsection{beacon:\{ip\}:\{ja4\}}

\begin{fullwidth}
\begin{table}[h]
\caption{\code{beacon:\{ip\}:\{ja4\}} key reference}
\label{tab:redis-beacon}
\begin{tabular}{@{}ll@{}}
\toprule
\textbf{Key pattern} & \code{beacon:\{ip\}:\{ja4\}} \\
\textbf{Type} & Sorted Set \\
\textbf{Score} & Unix timestamp in milliseconds \\
\textbf{Member} & Unix timestamp (same as score — member must be unique) \\
\textbf{TTL} & Managed by ZREMRANGEBYSCORE (30-minute sliding window) \\
\textbf{Writer} & Beaconing detector signal \\
\textbf{Reader} & Beaconing detector signal (ZRANGEBYSCORE) \\
\textbf{Example key} & \code{beacon:185.220.101.1:t13d1900...} \\
\bottomrule
\end{tabular}
\end{table}
\end{fullwidth}

\section{Connection State}
\index{Redis!connection state}

\subsection{conn:\{ip\}}

\begin{fullwidth}
\begin{table}[h]
\caption{\code{conn:\{ip\}} key reference — concurrent connections}
\label{tab:redis-conn}
\begin{tabular}{@{}ll@{}}
\toprule
\textbf{Key pattern} & \code{conn:\{ip\}} \\
\textbf{Type} & String (integer) \\
\textbf{TTL} & Managed by INCR/DECR lifecycle; expires after last connection closes \\
\textbf{Writer} & TCP connection handler (INCR on accept, DECR on close) \\
\textbf{Reader} & TCP behaviour signal \\
\textbf{Example} & \code{conn:185.220.101.1} → \code{"42"} \\
\bottomrule
\end{tabular}
\end{table}
\end{fullwidth}

\subsection{visit:\{ip\}}

\begin{fullwidth}
\begin{table}[h]
\caption{\code{visit:\{ip\}} key reference — return visitor tracking}
\label{tab:redis-visit}
\begin{tabular}{@{}ll@{}}
\toprule
\textbf{Key pattern} & \code{visit:\{ip\}} \\
\textbf{Type} & Hash \\
\textbf{Fields} & \code{first\_seen} (Unix timestamp), \code{total} (integer), \code{allowed} (integer), \code{blocked} (integer) \\
\textbf{TTL} & None (persistent for now; rotation policy TBD) \\
\textbf{Writer} & Action decider (HINCRBY on each action) \\
\textbf{Reader} & TCP behaviour signal \\
\bottomrule
\end{tabular}
\end{table}
\end{fullwidth}

\subsection{sess:\{ip\}}

\begin{fullwidth}
\begin{table}[h]
\caption{\code{sess:\{ip\}} key reference — session resumption ratio}
\label{tab:redis-sess}
\begin{tabular}{@{}ll@{}}
\toprule
\textbf{Key pattern} & \code{sess:\{ip\}} \\
\textbf{Type} & Hash \\
\textbf{Fields} & \code{total} (integer), \code{resumed} (integer) \\
\textbf{TTL} & None (persistent) \\
\textbf{Writer} & TLS handshake analyser \\
\textbf{Reader} & TCP behaviour signal (ratio = resumed/total) \\
\bottomrule
\end{tabular}
\end{table}
\end{fullwidth}

\section{HyperLogLog CIDR Density}
\index{Redis!HyperLogLog}
\index{HyperLogLog}

\subsection{hll:cidr24:\{cidr\} and hll:cidr48:\{cidr\}}

\begin{fullwidth}
\begin{table}[h]
\caption{HyperLogLog CIDR density keys}
\label{tab:redis-hll}
\begin{tabular}{@{}ll@{}}
\toprule
\textbf{Key patterns} & \code{hll:cidr24:\{cidr\}} (IPv4 /24), \code{hll:cidr48:\{cidr\}} (IPv6 /48) \\
\textbf{Type} & HyperLogLog \\
\textbf{Operations} & \code{PFADD} (add IP), \code{PFCOUNT} (count unique IPs) \\
\textbf{Memory} & Fixed 12KB per key, regardless of distinct IP count \\
\textbf{Error rate} & 0.81\% standard error \\
\textbf{Example key} & \code{hll:cidr24:185.220.101.0/24} \\
\textbf{Writer} & IP normalisation layer \\
\textbf{Reader} & Analytics node (CIDR density signals) \\
\textbf{TTL} & 30 days \\
\bottomrule
\end{tabular}
\end{table}
\end{fullwidth}

HyperLogLog is the only correct data structure for this use case. Tracking all distinct
IPs in a /24 in a Hash would require up to 256 entries per key; in a /48 it would be
unbounded. HyperLogLog provides an accurate-enough count with fixed memory cost.

\section{Enrichment Caches}
\index{Redis!enrichment cache}

\subsection{abuseipdb:\{ip\}}

\begin{fullwidth}
\begin{table}[h]
\caption{\code{abuseipdb:\{ip\}} key reference}
\label{tab:redis-abuseipdb}
\begin{tabular}{@{}ll@{}}
\toprule
\textbf{Key pattern} & \code{abuseipdb:\{ip\}} \\
\textbf{Type} & String \\
\textbf{Value} & JSON: \code{\{"score": 85, "cached\_at": 1712232896, "categories": [14, 21]\}} \\
\textbf{TTL} & \configkey{abuseipdb.cache\_ttl} seconds (default 3600) \\
\textbf{Writer} & AbuseIPDB signal (async, after API response) \\
\textbf{Reader} & AbuseIPDB signal (cache check) \\
\bottomrule
\end{tabular}
\end{table}
\end{fullwidth}

\subsection{rdap:\{ip\_or\_prefix\}}

\begin{fullwidth}
\begin{table}[h]
\caption{\code{rdap:\{ip\_or\_prefix\}} key reference}
\label{tab:redis-rdap}
\begin{tabular}{@{}ll@{}}
\toprule
\textbf{Key pattern} & \code{rdap:\{ip\_or\_prefix\}} \\
\textbf{Type} & String \\
\textbf{Value} & JSON: \code{\{"org": "Example ISP", "prefix": "185.220.0.0/16", "country": "DE"\}} \\
\textbf{TTL} & 86400 seconds (24h) \\
\textbf{Writer} & RDAP enrichment signal \\
\textbf{Reader} & RDAP enrichment signal (cache check) \\
\textbf{Example key} & \code{rdap:185.220.101.1}, \code{rdap:185.220.0.0/16} \\
\bottomrule
\end{tabular}
\end{table}
\end{fullwidth}

\subsection{geo:\{ip\}}

\begin{fullwidth}
\begin{table}[h]
\caption{\code{geo:\{ip\}} key reference — GeoIP country cache}
\label{tab:redis-geo}
\begin{tabular}{@{}ll@{}}
\toprule
\textbf{Key pattern} & \code{geo:\{ip\}} \\
\textbf{Type} & String \\
\textbf{Value} & ISO 3166-1 alpha-2 country code, e.g.\ \code{"RU"} \\
\textbf{TTL} & 86400 seconds (24h) \\
\textbf{Writer} & GeoIP lookup layer \\
\textbf{Reader} & GeoIP country block layer \\
\bottomrule
\end{tabular}
\end{table}
\end{fullwidth}

\section{Analytics Infrastructure}
\index{Redis!streams}
\index{Redis Streams}

\subsection{events — The Event Stream}

\begin{fullwidth}
\begin{table}[h]
\caption{\code{events} Redis Stream key reference}
\label{tab:redis-events}
\begin{tabular}{@{}ll@{}}
\toprule
\textbf{Key} & \code{events} \\
\textbf{Type} & Stream \\
\textbf{Writer} & Proxy (XADD, fire-and-forget) \\
\textbf{Reader} & Analytics node (XREADGROUP with consumer group \code{analytics}) \\
\textbf{Retention} & MAXLEN 100000 entries (approximate, using \code{XADD MAXLEN \textasciitilde}) \\
\textbf{Operations} & \code{XADD events * field1 val1 ...} (write), \code{XREADGROUP GROUP analytics consumer1 COUNT 100 STREAMS events >} (read) \\
\bottomrule
\end{tabular}
\end{table}
\end{fullwidth}

Every connection that completes the scoring pipeline produces one stream entry.
Bypassed connections (ALPN, whitelist) also produce a minimal entry for volume tracking.
The analytics node processes entries in batch using consumer groups, allowing multiple
analytics workers to process in parallel without duplicate processing.

\begin{notebox}[Why Streams not Pub/Sub]
Redis Pub/Sub is ephemeral — a message published while a subscriber is down is lost.
Redis Streams persist messages and allow consumer groups with acknowledgements. After
a restart or deployment, the analytics node replays missed events from where it left
off. This is critical for campaign detection, which requires historical context.
\end{notebox}

\subsection{findings:\{ip\}}

\begin{fullwidth}
\begin{table}[h]
\caption{\code{findings:\{ip\}} key reference — analytics findings feed-back}
\label{tab:redis-findings}
\begin{tabular}{@{}ll@{}}
\toprule
\textbf{Key pattern} & \code{findings:\{ip\}} \\
\textbf{Type} & String \\
\textbf{Value} & JSON: \code{\{"campaign": "scan-wave-14", "attribution\_score": 80, "coordinated": true\}} \\
\textbf{TTL} & 300 seconds \\
\textbf{Writer} & Analytics node \\
\textbf{Reader} & Pipeline layers 11 and 12 (attacker attribution, behavioural analysis) \\
\bottomrule
\end{tabular}
\end{table}
\end{fullwidth}

\section{Management and Audit}
\index{Redis!management}
\index{policy audit log}

\subsection{management:policy\_audit}

\begin{fullwidth}
\begin{table}[h]
\caption{\code{management:policy\_audit} key reference}
\label{tab:redis-audit}
\begin{tabular}{@{}ll@{}}
\toprule
\textbf{Key} & \code{management:policy\_audit} \\
\textbf{Type} & LIST \\
\textbf{Value} & JSON entries: \code{\{"timestamp": "...", "changed\_by": "admin@10.0.4.2", "item": "dial", "old\_value": 0, "new\_value": 50\}} \\
\textbf{TTL} & None (persistent; max 1000 entries via LTRIM) \\
\textbf{Writer} & Management UI (config changes); config reload process \\
\textbf{Reader} & Management UI audit log view \\
\textbf{Operations} & \code{LPUSH} (prepend new entry), \code{LTRIM 0 999} (keep last 1000) \\
\bottomrule
\end{tabular}
\end{table}
\end{fullwidth}

Every change to \configkey{security\_policy} is written here. Changes via the
Management UI are attributed to the operator's session IP. Changes via \code{SIGHUP}
config reload are attributed to \code{"config\_reload"}. This list provides the
complete audit trail for compliance purposes.

\section{Key Namespace Summary}
\index{Redis!key namespace}

\begin{fullwidth}
\begin{table}[h]
\caption{Complete Redis key namespace}
\label{tab:redis-namespace}
\begin{tabular}{@{}lllp{4.5cm}@{}}
\toprule
\textbf{Prefix} & \textbf{Type} & \textbf{TTL} & \textbf{Purpose} \\
\midrule
\code{ja4:blacklist} & SET & None & JA4 fingerprint blacklist \\
\code{ja4:whitelist} & SET & None & JA4 fingerprint whitelist \\
\code{ban:\{ip\}} & String & 300s–604800s & IP ban with escalating TTL \\
\code{rate:\{strategy\}:\{key\}} & Sorted Set & Managed & Sliding window rate limits \\
\code{beacon:\{ip\}:\{ja4\}} & Sorted Set & Managed & Beaconing IAT timestamps \\
\code{conn:\{ip\}} & String (int) & Managed & Concurrent connection count \\
\code{visit:\{ip\}} & Hash & None & Return visitor history \\
\code{sess:\{ip\}} & Hash & None & Session resumption ratio \\
\code{hll:cidr24:\{cidr\}} & HyperLogLog & 30d & Unique IPs per /24 (IPv4) \\
\code{hll:cidr48:\{cidr\}} & HyperLogLog & 30d & Unique IPs per /48 (IPv6) \\
\code{abuseipdb:\{ip\}} & String (JSON) & 3600s & AbuseIPDB score cache \\
\code{rdap:\{ip\}} & String (JSON) & 86400s & RDAP org lookup cache \\
\code{geo:\{ip\}} & String & 86400s & GeoIP country cache \\
\code{events} & Stream & MAXLEN 100k & Proxy event stream \\
\code{findings:\{ip\}} & String (JSON) & 300s & Analytics findings feed-back \\
\code{management:policy\_audit} & List & None (max 1000) & Policy change audit log \\
\bottomrule
\end{tabular}
\end{table}
\end{fullwidth}
